\section{Ethics}
\label{sec:ethics}

Following \citet{benotti-blackburn-2022-ethics}, I structure this section around research benefits, potential harms and vulnerable groups. Ethics statements in NLP papers have themselves become an object of study \citep{karamolegkou-etal-2025-ethical}, which supports treating them as substantive content. I cover the original paper first and then the issues that would arise if I published on my new data. The arXiv annotation was not completed, so the points about annotators below are prospective.

\paragraph{Research benefits.} REPaL needs only a relation definition and an unlabelled corpus, which could reduce the human annotation that RE otherwise requires. It also lowers the cost of building extractors for new relations. A replication adds value by reporting where results hold and where they do not. Here, that includes the mismatch between the paper's stated 10,000-instance pool and the released 100,000-instance pool.

\paragraph{Potential harms in the original paper.} The authors' ethics statement says that they do not verify the factualness of generated instances and ``simply count on the LLMs'', arguing that the text context is enough for sentence-level RE. They warn that adaptations needing factualness should be careful. I consider this a real limitation: seed examples in my logs name real historical figures, and those sentences are unverified. The paper also acknowledges that bias or errors can be introduced during synthesis. Its feedback stage targets the SLM's pattern biases, but that is not the same as addressing social bias. The evaluation data also contain sentences about named people, and FewRel relies on crowd workers' annotations.

\paragraph{Potential harms if I published on my new data.}
\emph{arXiv text reuse and author privacy.} Abstracts were collected through the official API within its rate limit. arXiv's terms permit use of descriptive metadata, defined to include titles and abstracts, under CC0 1.0 \citep{arxivtou}, but do not extend to storing and serving e-prints, and I collected no full text. Abstracts are not anonymous: the corpus keeps arXiv identifiers for attribution, and these link each sentence to its authors. Labels such as \texttt{outperforms}, attached to a named paper's work, could be read as evaluative claims about its authors. Text anonymisation methods such as differentially private rewriting \citep{meisenbacher-etal-2024-dp} exist, but would conflict with attribution and with the entity spans RE needs, so a release would instead need a takedown route for authors.

\emph{Annotator consent and risk.} Annotators would be participants in research, so the design would need assessment against the National Statement \citep{nhmrc-national-statement-2025}, including informed and voluntary consent, the freedom to withdraw, and proportionate risk. The risk appears low, as the task involves reading technical sentences, but the burden and any relationship between researcher and annotator would need attention. Because no annotators were recruited, there are no annotation data and no agreement figures, and I did not simulate any.

\emph{Responsible conduct and data management.} The Australian Code \citep{nhmrc-code-2018} frames responsible conduct and data management. In practice, I report shortfalls and non-clean runs, state the generator behind every result, and do not average across generators. Per-run configurations, logs and metrics are retained, and I keep a file manifest with checksums for the backup of results.

\emph{LLM-generated training data.} Synthetic examples can contain errors and encode the generator's biases, which the trained model can inherit. On SemEval, seed style differed clearly between GPT-4o and the Qwen models, showing that the generator can shape the data substantially. Local models also vary by version, so model identifiers are recorded for each run. Automated moral reasoning is an active research area \citep{pyatkin-etal-2023-clarifydelphi}, but I did not use such tools, and these judgements are my own.

\emph{Responsibility for AI-generated output.} Generative AI tools assisted with this project, as the unit permits. In line with Macquarie University's policy \citep{mq-ai-policy}, I remain responsible for this report and for checking AI-generated output.

\paragraph{Vulnerable groups.} Three groups deserve particular attention. People named in Wikipedia and news sentences cannot consent to automated extraction, and false relations could harm them. Authors, especially early-career researchers, could be affected by automated claims about their work. Prospective annotators could be vulnerable if recruited from students or others in a dependent relationship with the researcher, which is a further reason that consent would need to be genuinely voluntary.
